\honest{Rationality: Common Interest of Many Causes}{Eliezer Yudkowsky}{2009}

Many causes gain when rationality spreads, because it takes a little more rationality than usual to see their case: atheism, marijuana legalization, and, in an extreme case, my own institute, where I ``explicitly recursed on the job of creating rationalists.'' \nb{The institute is the Singularity Institute, now the Machine Intelligence Research Institute.} Not all the rationalists I create will join my project, and that is fine. ``You can't capture all the value you create.''

So causes could share standard material on rationality. Each would capture some of the value the others create. This requires that each be willing to ``shut up'' about its own cause, and not treat supporters as a fixed supply.

An unnamed person, taking over as president of an unnamed organization, pointed out that its message ``This is the best thing you can do'' had fared worse than the X Prize Foundation's ``Here is a cool thing you can do.'' \nb{This is one comparison, reported secondhand, between two different kinds of organization; the X Prize Foundation offers prizes for set, public goals.} People are not expected-utility maximizers but ``altruistic akrasics.'' Calling a cause the best adds little motivation, raises the burden of proof, and invites comparisons that sap motivation. \nb{A direct-mail experiment by Karlan and Wood found that adding scientific evidence of a charity's impact made no difference to giving overall. That is indirect support.} Arguing over the best cause ``may'' reduce the total supply of altruism. Dollars are fungible in theory, but they come from different mental accounts, and telling someone about another cause ``may increase the total amount they're willing to help.'' \nb{Field studies since then mostly agree. After deadly tornadoes, households gave more without giving less to other local charities (Deryugina and Marx, 2021). One study of major appeals found that they raised total giving while shifting the timing of gifts to other charities.}

There are limits. A project to find homes for stray puppies should ``probably'' be seen as exploiting ``bugs in human psychology for their personal gain.'' I give no evidence for that motive. The limit I mean is whether a project is a task ``you would wish to see done on behalf of humanity.''

So in general forums we should say ``Here is a cool rationalist project,'' and send anyone who asks which is best to a specialized subforum. Money flowing to the most underserved project will lower its returns anyway. Groups that will not share can be dropped from common resources.

If we all identified with a common ``Neo-Enlightenment,'' we would more often find ``ten of us in any given city.'' Rationality may not be ``the most important thing in the world,'' but it is ``a cool thing'' that many of us have in common.
